%(BEGIN_QUESTION)
% Copyright 2006, Tony R. Kuphaldt, released under the Creative Commons Attribution License (v 1.0)
% This means you may do almost anything with this work of mine, so long as you give me proper credit

En platina-RTD med $R_0$ = 1000 $\Omega$ og $\alpha$ = 0.00392 $\Omega$/$\Omega \cdot ^{o}$C varmes opp til 120$^{o}$ C. Beregn motstanden ved denne temperaturen.

\vskip 10pt

Beregn også temperaturen for samme RTD dersom målt motstand er 1043.8 $\Omega$.

\vskip 20pt \vbox{\hrule \hbox{\strut \vrule{} {\bf Forslag til sokratisk diskusjon} \vrule} \hrule}

\begin{itemize}
\item{} Forklar hvordan en RTD kan konstrueres med ulik basisresistans (f.eks. 1000 ohm kontra 100 ohm). Hvilke fysiske forhold i konstruksjonen bestemmer denne verdien?
\item{} Forklar hvorfor RTD-er ikke alle har samme alpha-verdi ($\alpha$).
\end{itemize}

\underbar{file i00406}
%(END_QUESTION)





%(BEGIN_ANSWER)

\noindent
{\bf Delvis svar:}

\vskip 10pt

$R_T$ = 1470.4 $\Omega$ ved 120$^{o}$ C

%(END_ANSWER)





%(BEGIN_NOTES)

$T$ = 11.173$^{o}$ C ved 1043.8 $\Omega$

\vfil \eject

\noindent
{\bf Oppsummeringsquiz:}

Beregn motstanden til en RTD med alpha = 0.00385 $\Omega$/$\Omega ^{o}$C og basisresistans 100 ohm (ved vannets frysepunkt), når temperaturen er 450 grader Fahrenheit.

\begin{itemize}
\item{} 189.4 $\Omega$
\vskip 5pt
\item{} 100.0 $\Omega$
\vskip 5pt
\item{} 260.9 $\Omega$
\vskip 5pt
\item{} 89.41 $\Omega$
\vskip 5pt
\item{} 196.3 $\Omega$
\vskip 5pt
\item{} 273.3 $\Omega$
\end{itemize}


%INDEX% Measurement, temperature: RTD

%(END_NOTES)
